\naslov{Večrazsežna normalna porazdelitev}

\begin{definicija}
  \pojem{Standardna večrazsežna normalna porazdelitev} je porazdelitev
  slučajnega vektorja $(Z_1, \ldots, Z_n)$, kjer so $Z_i \sim N(0, 1)$
  neodvisne.
\end{definicija}

\begin{opomba}
  To je $n$-razsežna porazdelitev z gostoto
  \[
	\phi_n(z_1, \ldots, z_n) = \inv{(2 \pi)^{n/2}} \exp\!\left( -
	  \frac{z_1^2}{2} - \cdots - \frac{z_n^2}{2} \right)
  \]
\end{opomba}

\begin{definicija}
  \pojem{Večrazsežna normalna porazdelitev} je poljubna porazdelitev slučajnega
  vektorja oblike $A \sv{Z} + \sv{\mu}$, kjer je $\sv{Z}$ slučajni večrazsežni
  normalni vektor, $A$ deterministična matrika, $\sv{\mu}$ pa deterministični
  vektor primerne velikosti.
\end{definicija}

\begin{opomba}
  Za vektor $\sv{X} = A \sv{Z} + \sv{\mu}$ je $E(\sv{X}) = \sv{\mu}$ in
  $\var(\sv{X}) = A A^T$.
\end{opomba}

\vprasanje{Definiraj večrazsežno normalno porazdelitev.}

Če je $\sv{Z}$ standardni $n$-razsežni normalni slučajni vektor, $A \in \R^{n
  \times n}$ deterministična matrika in $\mu \in \R^n$.
Tedaj ima $\sv{X} = A \sv{Z} + \sv{\mu}$ gostoto
\[
  f_{\sv{X}}(\sv{x})
  = \inv{(2 \pi)^{n/2} \sqrt{\det \Sigma}} \exp\!\left( - \pol (\sv{x} -
	\sv{\mu})^T \Sigma^{-1} (\sv{x} - \sv{\mu}) \right)
\]
za $\Sigma = A A^T$.

\begin{trditev}
  Če je $m \le n$, $\mu \in \R^m$ in ima $A \in \R^{m \times n}$ poln rang, je
  $A \sv{Z} + \sv{\mu}$ porazdeljena $N(\sv{\mu}, A A^T)$.
\end{trditev}

\begin{proof}
  Razcepimo lahko $A = B P Q$, kjer je $Q \in \R^{n \times n}$ ortogonalna, $P
  \in \R^{m \times n}$ koordinatna projekcija na prvih $m$ komponent, in $B \in
  \R^m$ neizrojena.
  Ker je $Q Q^T = I$, je $Q \sv{Z} \sim N(0, I)$, in so komponente tega vektorja
  neodvisne standardno normalne.
  Vektor $P Q \sv{Z}$ je prav tako standardni normalni z enakimi komponentami,
  le da jih ima $m$ namesto $n$.
  Torej je $A \sv{Z} + \sv{\mu} = B P Q \sv{Z} + \sv{\mu}$ porazdeljen $N(\mu, B
  B^T) = N(\mu, A A^T)$.
\end{proof}

\vprasanje{Kako je porazdeljen vektor $A \sv{Z} + \sv{\mu}$, kjer je $A$ polnega
  ranga?}

\begin{trditev}
  Porazdelitev poljubnega večrazsežnega normalnega slučajnega vektorja je
  natančno določena z njegovo pričakovano vrednostjo in kovariančno matriko.
\end{trditev}

\begin{proof}
  Če je $\sv{X}$ $m$-razsežen normalen slučajni vektor, po definiciji obstajata
  $A \in \R^{m \times n}$ in $\sv{\mu} \in \R^m$, da je $\sv{X} = A \sv{Z} +
  \sv{\mu}$ za $\sv{Z} \sim N(0, I_n)$.

  Naj bo $\tilde{\sv{Z}} \sim N(0, I_m)$ neodvisen od $\sv{Z}$.
  Definiramo $\sv{X}_k = A \sv{Z} + \inv{k} \tilde{\sv{Z}} + \sv{\mu}$.
  Iz izrekov Sluckega sledi $\sv{X}_k \xrightarrow[k \to \infty]{d} \sv{X}$.
  Ker sta $\sv{Z}$ in $\tilde{\sv{Z}}$ neodvisna, je vektor $(\sv{Z},
  \tilde{\sv{Z}}) \sim N(0, I_{n+m})$.
  Pišimo
  \[
	\sv{X}_k =
	\begin{bmatrix}
	  A & \inv{k} I_m
	\end{bmatrix}
	\cdot
	\begin{bmatrix}
	  \sv{Z} \\ \tilde{\sv{Z}}
	\end{bmatrix}
	+ \sv{\mu}.
  \]
  Ta matrika je polnega ranga, zato je porazdelitev slučajnega vektorja
  $\sv{X}_k$ odvisna le od pričakovane vrednosti $\sv{\mu}$ in kovariančne
  matrike $\var(\sv{X}_k) = A A^T + \inv{k^2} I_m$.
  Porazdelitev $\sv{X}$ je limita verjetnostnih porazdelitev $N(\sv{\mu}, A A^T
  + \inv{k^2} I_m)$; ta je natančno določena, torej je porazdelitev $\sv{X}$
  natančno določena z $\sv{\mu}$ in $A A^T$.
\end{proof}

\vprasanje{Dokaži: porazdelitev normalnega slučajnega vektorja je natančno
  določena s pričakovano vrednostjo in kovariančno matriko.}

\begin{trditev}
  Če je bločni slučajni vektor $(\sv{X}_1, \sv{X}_2)$ porazdeljen večrazsežno
  normalno in sta $\sv{X}_1$ ter $\sv{X}_2$ nekorelirana, sta tudi neodvisna.
\end{trditev}

\begin{proof}
  Velja
  \[
	(\sv{X}_1, \sv{X}_2) \sim N\!\left(
	  \begin{bmatrix}
		\sv{\mu}_1 \\ \sv{\mu}_2
	  \end{bmatrix},
	  \begin{bmatrix}
		\Sigma_1 & \\
		& \Sigma_2
	  \end{bmatrix}
	\right),
  \]
  to pa je tudi porazdelitev vektorja, kjer sta prvi in drugi blok nekorelirana.
\end{proof}

% LocalWords:  Sluckega
